%! Linear Combinations of Vectors — Unit 1, Lesson 1.1 — Warm-Up (Answer Key)
\documentclass[10pt]{article}
\usepackage{linalg-article}
\usepackage{linalg-key}   % pulls in -boxes; do NOT also load -boxes
\newcommand{\TallMath}[1]{$\displaystyle #1\rule[-1.4em]{0pt}{3.2em}$}
\newcommand{\vv}{\mathbf{v}}
\newcommand{\ww}{\mathbf{w}}

\begin{document}

\pageheader{Unit 1, Lesson 1.1}{Warm-Up --- Answer Key}
\namedateperiod

\vspace{0.12in}

\begin{notesbox}{Getting Ready: Skills We'll Use Today}
\small These are the Lesson~1.0 moves we combine today. Answer quickly.

\vspace{4pt}
\textbf{1. Scaling a vector.} Multiply every component by the scalar.
\begin{enumerate}[label=(\alph*), leftmargin=2.0em, itemsep=8pt]
  \item $3\begin{bmatrix} 2 \\ -1 \end{bmatrix} =
        \begin{bmatrix}{\color{keyred}\mathbf{6}}\\[2pt]{\color{keyred}\mathbf{-3}}\end{bmatrix}$
  \item $-2\begin{bmatrix} 1 \\ 4 \end{bmatrix} =
        \begin{bmatrix}{\color{keyred}\mathbf{-2}}\\[2pt]{\color{keyred}\mathbf{-8}}\end{bmatrix}$
\end{enumerate}

\vspace{8pt}
\textbf{2. Adding vectors.} Add matching components.
\[
  \begin{bmatrix} 2 \\ 1 \end{bmatrix} + \begin{bmatrix} 1 \\ 3 \end{bmatrix} =
  \begin{bmatrix}{\color{keyred}\mathbf{3}}\\[3pt]{\color{keyred}\mathbf{4}}\end{bmatrix}
\]

\vspace{6pt}
\textbf{3. Spotting a multiple.} Is $\begin{bmatrix} 4 \\ 2 \end{bmatrix}$ a scalar multiple of
$\begin{bmatrix} 2 \\ 1 \end{bmatrix}$? If so, by what number? Explain how you know.
\ansline{Yes --- it is $2\begin{bsmallmatrix} 2\\1 \end{bsmallmatrix}$. Both components scale
by the same factor $2$, so the two vectors point along the same line.}
\end{notesbox}

\begin{teachernote}
Item 3 is the pivot for today: ``same multiple of both components'' means the vectors are
\emph{parallel} (same line through the origin). That is exactly the case where two vectors
can \emph{not} reach the whole plane.
\end{teachernote}

\end{document}
